\section{AMC 12 12B 2006}
        \begin{enumerate}
        	\item (\textit{AMC 12 12B 2006}) What is $( - 1)^1 + ( - 1)^2 + \cdots + ( - 1)^{2006}$?


 $\text {(A) } - 2006 \qquad \text {(B) } - 1 \qquad \text {(C) } 0 \qquad \text {(D) } 1 \qquad \text {(E) } 2006$

 \textbf{Answer: }C



\textbf{Solution}

$(-1)^n=1$ if n is even and $-1$ if n is odd. So we have

 $-1+1-1+1-\cdots-1+1=0+0+\cdots+0+0=0 \Rightarrow \text{(C)}$



	\item (\textit{AMC 12 12B 2006}) For real numbers $x$ and $y$, define $x\spadesuit y = (x + y)(x - y)$.  What is $3\spadesuit(4\spadesuit 5)$?


 $\text {(A) } - 72 \qquad \text {(B) } - 27 \qquad \text {(C) } - 24 \qquad \text {(D) } 24 \qquad \text {(E) } 72$

 \textbf{Answer: }A



\textbf{Solution}

$4\spadesuit 5=-9$
$3\spadesuit -9=-72 \Rightarrow \text{(A)}$



	\item (\textit{AMC 12 12B 2006}) A football game was played between two teams, the Cougars and the Panthers.  The two teams scored a total of 34 points, and the Cougars won by a margin of 14 points.  How many points did the Panthers score?


 $\text {(A) } 10 \qquad \text {(B) } 14 \qquad \text {(C) } 17 \qquad \text {(D) } 20 \qquad \text {(E) } 24$

 \textbf{Answer: }A



\textbf{Solution 1}

If the Cougars won by a margin of 14 points, then the Panthers' score would be half of (34-14). That's 10 $\Rightarrow \boxed{\text{(A)}}$.

 



\textbf{Solution 2}

Let the Panthers' score be $x$. The Cougars then scored $x+14$. Since the teams combined scored $34$, we get $x+x+14=34 \\ \rightarrow 2x+14=34 \\ \rightarrow 2x=20 \\ \rightarrow x = 10$, 

 and the answer is  $\boxed{\text{(A)}}$.

 


	\item (\textit{AMC 12 12B 2006}) Mary is about to pay for five items at the grocery store.  The prices of the items are $\textdollar7.99$, $\textdollar4.99$, $\textdollar2.99$, $\textdollar1.99$, and $\textdollar0.99$.  Mary will pay with a twenty-dollar bill.  Which of the following is closest to the percentage of the $\textdollar20.00$ that she will receive in change?


 $\text {(A) } 5 \qquad \text {(B) } 10 \qquad \text {(C) } 15 \qquad \text {(D) } 20 \qquad \text {(E) } 25$

 \textbf{Answer: }A



\textbf{Solution}

The total price of the items is $(8-.01)+(5-.01)+(3-.01)+(2-.01)+(1-.01)=19-.05=18.95$
$20-18.95=1.05$
$\frac{1.05}{20}=.0525 \Rightarrow \text{(A)}$

 \textbf{Alternative Solution} We can round the prices to $8$, $5$, $3$, $2$, and $1$. 

 So $20 - (8+5+3+2+1) = 1$

 We can make an equation: $20* \frac{x}{100}=1$

 If we simplify the equation to "x", we get $x=5$



	\item (\textit{AMC 12 12B 2006}) John is walking east at a speed of 3 miles per hour, while Bob is also walking east, but at a speed of 5 miles per hour.  If Bob is now 1 mile west of John, how many minutes will it take for Bob to catch up to John?


 $\text {(A) } 30 \qquad \text {(B) } 50 \qquad \text {(C) } 60 \qquad \text {(D) } 90 \qquad \text {(E) } 120$

 \textbf{Answer: }A



\textbf{Solution}

The speed that Bob is catching up to John is $5-3=2$ miles per hour. Since Bob is one mile behind John, it will take $\frac{1}{2} \Rightarrow \text{(A)}$ of an hour to catch up to John.

 


	\item (\textit{AMC 12 12B 2006}) Francesca uses 100 grams of lemon juice, 100 grams of sugar, and 400 grams of water to make lemonade.  There are 25 calories in 100 grams of lemon juice and 386 calories in 100 grams of sugar.  Water contains no calories.  How many calories are in 200 grams of her lemonade?


 $\text {(A) } 129 \qquad \text {(B) } 137 \qquad \text {(C) } 174 \qquad \text {(D) } 223 \qquad \text {(E) } 411$

 \textbf{Answer: }B



\textbf{Solution}

Francesca makes a total of $100+100+400=600$ grams of lemonade, and in those $600$ grams, there are $25$ calories from the lemon juice and $386$ calories from the sugar, for a total of $25+386=411$ calories per $600$ grams. We want to know how many calories there are in $200=600/3$ grams, so we just divide $411$ by $3$ to get $137\implies\boxed{(\text{B})}$.

 


	\item (\textit{AMC 12 12B 2006}) Mr. and Mrs. Lopez have two children.  When they get into their family car, two people sit in the front, and the other two sit in the back.  Either Mr. Lopez or Mrs. Lopez must sit in the driver's seat.  How many seating arrangements are possible?


 $\text {(A) } 4 \qquad \text {(B) } 12 \qquad \text {(C) } 16 \qquad \text {(D) } 24 \qquad \text {(E) } 48$

 \textbf{Answer: }B



\textbf{Solution 1}

First, we seat the children.

 The first child can be seated in $3$ spaces.

 The second child can be seated in $2$ spaces.

 Now there are $2 \times 1$ ways to seat the adults.

 $3 \times 2 \times 2 = 12 \Rightarrow \text{(B)}$




\textbf{Solution 2}

There are only two possible occupants for the driver's seat. After the driver is chosen, any of the remaining three people can sit in the front, and there are two arrangements for the other two people in the back. Thus, there are $2\cdot 3\cdot 2 = \boxed{12}$ possible seating arrangements. ~ aopsav (Credit to AoPS Alcumus)

 Alternative solution:

 If there was no restriction, there would be 4!=24 ways to sit. However, only 2/4 of the people can sit in the driver's seat, so our answer is $\frac{2}{4}\cdot 24= 12 \Rightarrow \text{(B)}$



	\item (\textit{AMC 12 12B 2006}) The lines $x = \frac 14y + a$ and $y = \frac 14x + b$ intersect at the point $(1,2)$.  What is $a + b$?


 $\text {(A) } 0 \qquad \text {(B) } \frac 34 \qquad \text {(C) } 1 \qquad \text {(D) } 2 \qquad \text {(E) } \frac 94$

 \textbf{Answer: }E



\textbf{Solution 1}

$4x-4a=y$
$4x-4a=\frac{1}{4}x+b$
$4\cdot1-4a=\frac{1}{4}\cdot1+b=2$
$a=\frac{1}{2}$
$b=\frac{7}{4}$
$a+b=\frac{9}{4} \Rightarrow \fbox{(E)}$




\textbf{Solution 2}

Add both equations:

 $\begin{cases}x=\frac{1}{4}y+a\\ y=\frac{1}{4}x+b \end{cases}$

 Simplify:

 $\frac{4}{4}(x+y)=\frac{1}{4}(x+y)+(a+b)$

 Isolate our solution:

 $\frac{3}{4}(x+y)=a+b$

 Substitute the point of intersection $[x=1, y=2]$
$a+b=\frac{3}{4}\cdot3=\frac{9}{4} \Rightarrow \fbox{(E)}$




\textbf{Solution 3}

Plugging in $(1,2)$ into the first equation, and solving for $a$ we get $a$ as $\frac{1}{2}$.

 Doing the same for the second equation for the second equation, we get $b$ as $\frac{7}{4}$

 Adding $a+b = \frac{1}{2} + \frac{7}{4} = \frac{9}{4} \Rightarrow \fbox{(E)}$



	\item (\textit{AMC 12 12B 2006}) How many even three-digit integers have the property that their digits, read left to right, are in strictly increasing order? 


 $\text {(A) } 21 \qquad \text {(B) } 34 \qquad \text {(C) } 51 \qquad \text {(D) } 72 \qquad \text {(E) } 150$

 \textbf{Answer: }B



\textbf{Solution 1 (Alcumus Edition)}

Let the integer have digits $a$, $b$, and $c$, read left to right. Because $1 \leq a<b<c$, none of the digits can be zero and $c$ cannot be 2. If $c=4$, then $a$ and $b$ must each be chosen from the digits 1, 2, and 3. Therefore there are $\binom{3}{2}=3$ choices for $a$ and $b$, and for each choice there is one acceptable order. Similarly, for $c=6$ and $c=8$ there are, respectively, $\binom{5}{2}=10$ and $\binom{7}{2}=21$ choices for $a$ and $b$. Thus there are altogether $3+10+21=\boxed{34}$ such integers.

 


 (Edited by HMSSONI82)

 



\textbf{Solution 2}

Let's set the middle (tens) digit first. The middle digit can be anything from 2-7 (If it was 1 we would have the hundreds digit to be 0, if it was more than 7, the ones digit cannot be even).

 If it was 2, there is 1 possibility for the hundreds digit, 3 for the ones digit.If it was 3, there are 2 possibilities for the hundreds digit, 3 for the ones digit.If it was 4, there are 3 possibilities for the hundreds digit, and 2 for the ones digit,

 and so on.

 So, the answer is $3(1+2)+2(3+4)+1(5+6)=\boxed{34} \Rightarrow B$.

 



\textbf{Solution 3}

The last digit is 4, 6, or 8.

 If the last digit is $x$, the possibilities for the first two digits correspond to 2-element subsets of $\{1,2,\dots,x-1\}$.

 Thus the answer is ${3\choose 2} + {5\choose 2} + {7\choose 2} = 3 + 10 + 21 = \boxed{34}$.

 



\textbf{Solution 4}

The answer must be half of a triangular number (evens and decreasing/increasing) so $\boxed{34}$ or the letter B.-

 



\textbf{Solution 5 (Forward Casework + Listing)}

Casework:

 For the sake of simplicity, we are going to call the number $\overline{abc}$.

 1. If $a=1$:

 a. $c=2$. No such number exists.

 b. $c=4$. $b=2, 3$.

 c. $c=6$. $b=2, 3, 4, 5$.

 d. $c=8$. $b=2, 3, 4, 5, 6, 7$.

 2. If $a=2$: continue as above.

 We can count up that there are 34 such integers, so select $\boxed{B}$.

 


	\item (\textit{AMC 12 12B 2006}) In a triangle with integer side lengths, one side is three times as long as a second side, and the length of the third side is 15. What is the greatest possible perimeter of the triangle?


 $\text {(A) } 43 \qquad \text {(B) } 44 \qquad \text {(C) } 45 \qquad \text {(D) } 46 \qquad \text {(E) } 47$

 \textbf{Answer: }A



\textbf{Solution}

If the second size has length x, then the first side has length 3x, and we have the third side which has length 15. By the triangle inequality, we have: 
 \[\\ x+15>3x \Rightarrow 2x<15 \Rightarrow x<7.5\] Now, since we want the greatest perimeter, we want the greatest integer x, and if $x<7.5$ then $x=7$. Then, the first side has length $3*7=21$, the second side has length $7$, the third side has length $15$, and so the perimeter is $21+7+15=43 \Rightarrow \boxed{\text {(A)}}$.

 


	\item (\textit{AMC 12 12B 2006}) Joe and JoAnn each bought 12 ounces of coffee in a 16-ounce cup.  Joe drank 2 ounces of his coffee and then added 2 ounces of cream.  JoAnn added 2 ounces of cream, stirred the coffee well, and then drank 2 ounces.  What is the resulting ratio of the amount of cream in Joe's coffee to that in JoAnn's coffee?


 $\text {(A) } \frac 67 \qquad \text {(B) } \frac {13}{14} \qquad \text {(C) } 1 \qquad \text {(D) } \frac {14}{13} \qquad \text {(E) } \frac 76$

 \textbf{Answer: }E



\textbf{Solution}

Joe has 2 ounces of cream, as stated in the problem.

 JoAnn had 14 ounces of liquid, and drank $\frac{1}{7}$ of it. Therefore, she drank $\frac{1}{7}$ of her cream, leaving her $2*\frac{6}{7}$.

 $\frac{2}{2*\frac{6}{7}}=\frac{7}{6} \Rightarrow \boxed{\text{(E)}}$



	\item (\textit{AMC 12 12B 2006}) The parabola $y=ax^2+bx+c$ has vertex $(p,p)$ and $y$-intercept $(0,-p)$, where $p\ne 0$. What is $b$?


 $\text {(A) } -p \qquad \text {(B) } 0 \qquad \text {(C) } 2 \qquad \text {(D) } 4 \qquad \text {(E) } p$

 \textbf{Answer: }D



\textbf{Solution 1}

Substituting $(0,-p)$, we find that $y = -p = a(0)^2 + b(0) + c = c$, so our parabola is $y = ax^2 + bx - p$. 

 The x-coordinate of the vertex of a parabola is given by $x = p = \frac{-b}{2a} \Longleftrightarrow a = \frac{-b}{2p}$. Additionally, substituting $(p,p)$, we find that $y = p = a(p)^2 + b(p) - p \Longleftrightarrow ap^2 + (b-2)p = \left(\frac{-b}{2p}\right)p^2 + (b-2)p = p\left(\frac b2-2\right) = 0$. Since it is given that $p \neq 0$, then $\frac{b}{2} = 2 \Longrightarrow b = 4\ \mathrm{(D)}$.

 



\textbf{Solution 2}

A parabola with the given equation and with vertex $(p,p)$ must have equation $y=a(x-p)^2+p$. Because the $y$-intercept is $(0,-p)$ and $p\ne 0$, it follows that $a=-2/p$. Thus
 \[y=-\frac{2}{p}(x^2-2px+p^2)+p=-\frac{2}{p}x^2+4x-p,\] so $\boxed{b=4}$.

 


	\item (\textit{AMC 12 12B 2006}) Rhombus $ABCD$ is similar to rhombus $BFDE$. The area of rhombus $ABCD$ is 24, and $\angle BAD = 60^\circ$. What is the area of rhombus $BFDE$?


 
                    \begin{figure}[H]
                        \centering
                        \includegraphics[width=0.5\textwidth]{Problem/AMC_12_12B_2006_13_Problem_Figure_1.png}
                    \end{figure}
                    
$\textrm{(A) } 6 \qquad \textrm{(B) } 4\sqrt {3} \qquad \textrm{(C) } 8 \qquad \textrm{(D) } 9 \qquad \textrm{(E) } 6\sqrt {3}$

 \textbf{Answer: }C



\textbf{Solution 1}

The ratio of any length on $ABCD$ to a corresponding length on $BFDE^2$ is equal to the ratio of their areas. Since $\angle BAD=60$, $\triangle ADB$ and $\triangle DBC$ are equilateral. $DB$, which is equal to $AB$, is the diagonal of rhombus $ABCD$. Therefore, $AC=\frac{DB(2)}{2\sqrt{3}}=\frac{DB}{\sqrt{3}}$. $DB$ and $AC$ are the longer diagonal of rhombuses $BEDF$ and $ABCD$, respectively. So the ratio of their areas is $(\frac{1}{\sqrt{3}})^2$ or $\frac{1}{3}$. One-third the area of $ABCD$ is equal to $8$. So the answer is $\boxed{\text{C}}$.

 



\textbf{Solution 2}

Draw the line $\overline{DB}$ to form an equilateral triangle, since $\angle BAD=60$, and line segments $\overline{AB}$ and $\overline{AD}$ are equal in length. To find the area of the smaller rhombus, we only need to find the value of any arbitrary base, then square the result. To find the value of the base, use the line we just drew and connect it to point $E$ at a right angle along line $\overline{DB}$. Call the connected point $P$, with triangles $DPE$ and $EPB$ being 30-60-90 triangles, meaning we can find the length of $\overline{ED}$ or $\overline{EB}$. The base of $ABCD$ must be $\sqrt{24}$, and half of that length must be $\sqrt{6}$(triangles $DPE$ and $EPB$ are congruent by $SSS$). Solving for the third length yields $\sqrt{8}$, which we square to get the answer $\boxed{\text{C}}$.

 



\textbf{Solution 3}

Draw line DB, cutting rhombus BFDE into two triangles which fit nicely into 2/3 of equilateral triangle ABD. Thus the area of BFDE is (2/3)*12=8.

 


	\item (\textit{AMC 12 12B 2006}) Elmo makes $N$ sandwiches for a fundraiser. For each sandwich he uses $B$ globs of peanut butter at $4$ cents per glob and $J$ blobs of jam at $5$ cents per blob. The cost of the peanut butter and jam to make all the sandwiches is $\textdollar 2.53$. Assume that $B$, $J$ and $N$ are all positive integers with $N>1$. What is the cost of the jam Elmo uses to make the sandwiches? 


 $\mathrm{(A)}\ 1.05 \qquad \mathrm{(B)}\ 1.25 \qquad \mathrm{(C)}\ 1.45 \qquad \mathrm{(D)}\ 1.65 \qquad \mathrm{(E)}\ 1.85$

 \textbf{Answer: }D



\textbf{Solution}

From the given, we know that 

 $253=N(4B+5J)$ (The numbers are in cents)

 since $253=11\cdot23$, and since $N$ is an integer, then $4B+5J=11$ or $23$. It is easily deduced that $11$ is impossible to make with $B$ and $J$ integers, so $N=11$ and $4B+5J=23$. Then, it can be guessed and checked quite simply that if $B=2$ and $J=3$, then $4B+5J=4(2)+5(3)=23$. The problem asks for the total cost of jam, or $N(5J)=11(15)=165$ cents, or $1.65\implies\mathrm{(D)}$



	\item (\textit{AMC 12 12B 2006}) Circles with centers $O$ and $P$ have radii 2 and 4, respectively, and are externally tangent.  Points $A$ and $B$ are on the circle centered at $O$, and points $C$ and $D$ are on the circle centered at $P$, such that $\overline{AD}$ and $\overline{BC}$ are common external tangents to the circles.  What is the area of hexagon $AOBCPD$?


 
                    \begin{figure}[H]
                        \centering
                        \includegraphics[width=0.5\textwidth]{Problem/AMC_12_12B_2006_15_Problem_Figure_1.png}
                    \end{figure}
                    
$\textbf{(A) } 18\sqrt {3} \qquad \textbf{(B) } 24\sqrt {2} \qquad \textbf{(C) } 36 \qquad \textbf{(D) } 24\sqrt {3} \qquad \textbf{(E) } 32\sqrt {2}$

 \textbf{Answer: }B



\textbf{Solution}

Draw the altitude from $O$ onto $DP$ and call the point $H$. Because $\angle OAD$ and $\angle ADP$ are right angles due to being tangent to the circles, and the altitude creates $\angle OHD$ as a right angle. $ADHO$ is a rectangle with $OH$ bisecting $DP$. The length $OP$ is $4+2$ and $HP$ has a length of $2$, so by pythagorean's, $OH$ is $\sqrt{32}$.

 $2 \cdot \sqrt{32} + \frac{1}{2}\cdot2\cdot \sqrt{32} = 3\sqrt{32} = 12\sqrt{2}$, which is half the area of the hexagon, so the area of the entire hexagon is $2\cdot12\sqrt{2} = \boxed{(B)24\sqrt{2}}$




\textbf{Solution 2}

$ADPO$ and $OPBC$ are congruent right trapezoids with legs $2$ and $4$ and with $OP$ equal to $6$. Draw an altitude from $O$ to either $DP$ or $CP$, creating a rectangle with width $2$ and base $x$, and a right triangle with one leg $2$, the hypotenuse $6$, and the other $x$. Usingthe \href{https://artofproblemsolving.com/wiki/index.php?title=Pythagorean\_theorem}{Pythagorean theorem}, $x$ is equal to $4\sqrt{2}$, and $x$ is also equal to the height of the trapezoid. The area of the trapezoid is thus $\frac{1}{2}\cdot(4+2)\cdot4\sqrt{2} = 12\sqrt{2}$, and the total area is two trapezoids, or $\boxed{24\sqrt{2}}$.

 



\textbf{Solution 3}

We can create two congruent rectangles and two congruent right triangles in the figure. We see that each rectangle will have a side length of $2$, and each triangle will therefore also have a base of $2$.

 The hypotenuse of each right triangle has a length of $6$, so we use the Pythagorean theorem and find the length of the height of each triangle as follows: 
 \[\sqrt{6^2 - 2^2} = \sqrt{36 - 4} = 4\sqrt{2}\]

 The area of each individual triangle is $\frac{2 \cdot 4\sqrt{2}}{2} = 4\sqrt{2}$. We multiply this by two to get the area of both triangles combined, which gives $8\sqrt{2}$.

 Now to find the area of each rectangle, we simply multiply $2$ by $4\sqrt{2}$ and double this, which gives the combined area of the two rectangles to be $16\sqrt{2}$. We add $8\sqrt{2}$ to this to give our answer of $\boxed{\textbf{(B) } 24\sqrt{2}}$.

 


	\item (\textit{AMC 12 12B 2006}) Regular hexagon $ABCDEF$ has vertices $A$ and $C$ at $(0,0)$ and $(7,1)$, respectively. What is its area?


 $\mathrm{(A)}\ 20\sqrt {3} \qquad \mathrm{(B)}\ 22\sqrt {3} \qquad \mathrm{(C)}\ 25\sqrt {3} \qquad \mathrm{(D)}\ 27\sqrt {3} \qquad \mathrm{(E)}\ 50$

 \textbf{Answer: }C



\textbf{Solution}

To find the area of the regular hexagon, we only need to calculate the side length.a distance of $\sqrt{7^2+1^2} = \sqrt{50} = 5\sqrt{2}$ apart.  Half of this distance is the length of the longer leg of the right triangles.  Therefore, the side length of the hexagon is $\frac{5\sqrt{2}}{2}\cdot\frac{1}{\sqrt{3}}\cdot2 = \frac{5\sqrt{6}}{3}$.

 The apothem is thus $\frac{1}{2}\cdot\frac{5\sqrt{6}}{3}\cdot\sqrt{3} = \frac{5\sqrt{2}}{2}$, yielding an area of $\frac{1}{2}\cdot10\sqrt{6}\cdot\frac{5\sqrt{2}}{2}=25\sqrt{3} \implies \mathrm{(C)}$.

 



\textbf{Solution 2}

Solution 2 has the exact same solution as Solution 1 denoted, but instead, we do not need to know the value of the apothem. We could just apply s, which is the side length in this problem, $\frac{5\sqrt{6}}{3}$ into the hexagon area formula, $\frac{3  (\frac{5\sqrt{6}}{3})^2 \sqrt{3} }{2}=25\sqrt{3} \implies \mathrm{(C)}$



	\item (\textit{AMC 12 12B 2006}) For a particular peculiar pair of dice, the probabilities of rolling $1$, $2$, $3$, $4$, $5$ and $6$ on each die are in the ratio $1:2:3:4:5:6$. What is the probability of rolling a total of $7$ on the two dice?


 $\mathrm{(A)}\ \frac 4{63} \qquad \mathrm{(B)}\ \frac 18  \qquad \mathrm{(C)}\ \frac 8{63} \qquad \mathrm{(D)}\ \frac 16 \qquad \mathrm{(E)}\ \frac 27$

 \textbf{Answer: }C



\textbf{Solution}

The probability of getting an $x$ on one of these dice is $\frac{x}{21}$.

 The probability of getting $1$ on the first and $6$ on the second die is $\frac 1{21}\cdot\frac 6{21}$. Similarly we can express the probabilities for the other five ways how we can get a total $7$. (Note that we only need the first three, the other three are symmetric.)

 Summing these, the probability of getting a total $7$ is:
 \[2\cdot\left( \frac 1{21}\cdot\frac 6{21} + \frac 2{21}\cdot\frac 5{21} + \frac 3{21}\cdot\frac 4{21} \right) =  \frac{56}{441} = \boxed{\frac{8}{63}}\]

 See also \href{https://artofproblemsolving.com/wiki/index.php?title=2016\_AIME\_I\_Problems/Problem\_2}{2016 AIME I Problems/Problem 2}

 


	\item (\textit{AMC 12 12B 2006}) An object in the plane moves from one lattice point to another. At each step, the object may move one unit to the right, one unit to the left, one unit up, or one unit down. If the object starts at the origin and takes a ten-step path, how many different points could be the final point?


 $\mathrm{(A)}\ 120 \qquad \mathrm{(B)}\ 121 \qquad \mathrm{(C)}\ 221 \qquad \mathrm{(D)}\ 230 \qquad \mathrm{(E)}\ 231$

 \textbf{Answer: }B



\textbf{Solution 1}

Let the starting point be $(0,0)$. After $10$ steps we can only be in locations $(x,y)$ where $|x|+|y|\leq 10$. Additionally, each step changes the \href{https://artofproblemsolving.com/wiki/index.php?title=Parity}{parity} of exactly one coordinate. Hence after $10$ steps we can only be in locations $(x,y)$ where $x+y$ is even. It can easily be shown that each location that satisfies these two conditions is indeed reachable.

 Once we pick $x\in\{-10,\dots,10\}$, we have $11-|x|$ valid choices for $y$, giving a total of $\boxed{121}$ possible positions.

 



\textbf{Solution 2}

$10$ moves results in a lot of possible endpoints, so we try small cases first.

 If the object only makes $1$ move, it is obvious that there are only 4 possible points that the object can move to.If the object makes $2$ moves, it can move to $(0, 2)$, $(1, 1)$, $(2, 0)$, $(1, -1)$, $(0, -2)$, $(-1, -1)$, $(-2, 0)$ as well as $(0, 0)$, for a total of $9$ moves.If the object makes $3$ moves, it can end up at $(0, 3)$, $(2, 1)$, $(1, 2)$, $(3, 0)$, $(2, -1)$, $(1, -2)$, $(0, -3)$, $(-1, -2)$, $(-2, -1)$, $(0, -3)$ etc. for a total of $16$ moves.

 At this point we can guess that for n moves, there are $(n + 1)^2$ different endpoints. Thus, for 10 moves, there are $11^2 = 121$ endpoints, and the answer is $\boxed{B}$.

 


	\item (\textit{AMC 12 12B 2006}) Mr. Jones has eight children of different ages. On a family trip his oldest child, who is 9, spots a license plate with a 4-digit number in which each of two digits appears two times. "Look, daddy!" she exclaims. "That number is evenly divisible by the age of each of us kids!" "That's right," replies Mr. Jones, "and the last two digits just happen to be my age." Which of the following is not the age of one of Mr. Jones's children?


 $\mathrm{(A)}\ 4 \qquad \mathrm{(B)}\ 5 \qquad \mathrm{(C)}\ 6 \qquad \mathrm{(D)}\ 7 \qquad \mathrm{(E)}\ 8$

 \textbf{Answer: }B



\textbf{Solution}

First, The number of the plate is divisible by $9$ and in the form of$aabb$, $abba$ or $abab$. 

 We can conclude straight away that $a+b= 9$ using the $9$ divisibility rule. 

 If $b=1$, the number is not divisible by $2$ (unless it's $1818$, which is not divisible by $4$), which means there are no $2$, $4$, $6$, or $8$ year olds on the car, but that can't be true, as that would mean there are less than $8$ kids on the car. 

 If $b=2$, then the only possible number is $7272$. $7272$ is divisible by $4$, $6$, and $8$, but not by $5$ and $7$, so that doesn't work. 

 If $b=3$, then the only number is $6336$, also not divisible by $5$ or $7$. 

 If $b=4$, the only number is $5544$. It is divisible by $4$, $6$, $7$, and $8$. 

 Therefore, we conclude that the answer is $\mathrm{(B)}\ 5$

 \textbf{NOTE}: Automatically, since there are 8 children and all of their ages are less than or equal to 9 and are different, the answer choices can be narrowed down to $5$ or $8$.

 \textbf{Alternate Solution} We know that the number of the plate is divisible by $9$ and in the form $aabb$, $abba$, or $abab$ for distinct digits $a,b$.

 Using the divisibility rule for $9$, we can conclude that $a+b\equiv 0\pmod{9}$.

 We also know that the number of the plate is even, because you can only discard one number from the integers $1$ through $9$, inclusive ($8$ children, oldest is $9$), and there's always going to be an even number left.

 If one of the children was $5$ years old, then the plate number would be divisible by both $5$ and $2$. Thus, the units digit must be $0$. Then, the possible form of the plate number would be $aa00$, $0bb0$, or $a0a0$. The first case is not possible because $00$ is not a possible age for the father.

 We have the remaining forms $0bb0$ and $a0a0$.

 $\textbf{Case 1: 0bb0:}$We know that $2b\equiv 0\pmod{9}$, so $b$ has to be either $0$ or $9$. $b$ can't be $0$ because $a, b$ are distinct. So, $b=9$. However, the number $0990$ is not divisible by both $4$ and $7$, so whichever number you discard, the other one will still be there.  This creates a contradiction, so Case $1$ cannot be true.

 $\textbf{Case 2: a0a0:}$Similarly to Case $1$, we can determine that $a$ has to be $9$. However, the number $9090$ is not divisible by both $4$ and $7$. Using the same logic in case $1$, we conclude that Case $2$ cannot be true.

 Since we have disproven all of our cases, we know that it is impossible for one of the children to be $\mathrm{(B)}\ 5$ years old.

 \textbf{NOTE}: This might look tedious, but it only took me around 30 seconds to do on paper.

 


	\item (\textit{AMC 12 12B 2006}) Let $x$ be chosen at random from the interval $(0,1)$. What is the probability that
$\lfloor\log_{10}4x\rfloor - \lfloor\log_{10}x\rfloor = 0$?
Here $\lfloor x\rfloor$ denotes the greatest integer that is less than or equal to $x$.


 

$\mathrm{(A)}\ \frac 18 \qquad \mathrm{(B)}\ \frac 3{20} \qquad \mathrm{(C)}\ \frac 16 \qquad \mathrm{(D)}\ \frac 15  \qquad \mathrm{(E)}\ \frac 14$

 \textbf{Answer: }C



\textbf{Solution}

Let $k$ be an arbitrary integer. For which $x$ do we have $\lfloor\log_{10}4x\rfloor = \lfloor\log_{10}x\rfloor = k$?

 The equation $\lfloor\log_{10}x\rfloor = k$ can be rewritten as $10^k \leq x < 10^{k+1}$. The second one gives us $10^k \leq 4x < 10^{k+1}$. Combining these, we get that both hold at the same time if and only if $10^k \leq x < \frac{10^{k+1}}4$.

 Hence for each integer $k$ we get an interval of values for which $\lfloor\log_{10}4x\rfloor - \lfloor\log_{10}x\rfloor = 0$. These intervals are obviously pairwise disjoint.

 For any $k\geq 0$ the corresponding interval is disjoint with $(0,1)$, so it does not contribute to our answer. On the other hand, for any $k<0$ the entire interval is inside $(0,1)$. Hence our answer is the sum of the lengths of the intervals for $k<0$.

 For a fixed $k$ the length of the interval $\left[ 10^k, \frac{10^{k+1}}4 \right)$ is $\frac 32\cdot 10^k$. 

 This means that our result is $\frac 32 \left( 10^{-1} + 10^{-2} + \cdots \right) = \frac 32 \cdot \frac 19 = \boxed{\frac 16}$.

 



\textbf{Solution 2}

The largest value for $x$ is $10^{0}$. If $x > 10^{-1}$, then $\lfloor\log_{10}4x\rfloor$ doesn't fulfill the condition unless $10^{-2} \leq x < 0.25 * 10^{-1}$. The same holds when you get smaller, because $x = 0.25 * 10^{n}$ for $n \leq 0$ is the lowest value such that $4x$ becomes a higher power of 10.

 Recognize that this is a geometric sequence. The probability of choosing $x$ such that  $\lfloor\log_{10}4x\rfloor$ and $\lfloor\log_{10}x\rfloor$ both equal $-1$ is $(9/10)* (15/90) =15/100$, because there is a 90 percent chance of choosing $x > 10^{-1}$, and only values of $x$ between $10^{-1}$ and $0.25*10^{0}$ work in this case. Then, for $x$ such that $\lfloor\log_{10}4x\rfloor$ and $\lfloor\log_{10}x\rfloor$ both equal $-2$, you have $(1/10) * ((9/10) *(15/90))$. This is a geometric series with ratio $1/10$. Using $a/(1-r)$ for the sum of an infinite geometric sequence, we get $(15/100)/(1-(1/10)) =  \boxed{\frac 16}$.

 Solution by Halt\_CatchFire

 


	\item (\textit{AMC 12 12B 2006}) Rectangle $ABCD$ has area $2006$.  An ellipse with area $2006\pi$ passes through $A$ and $C$ and has foci at $B$ and $D$.  What is the perimeter of the rectangle? (The area of an ellipse is $ab\pi$ where $2a$ and $2b$ are the lengths of the axes.)


 $\mathrm{(A)}\ \frac {16\sqrt {2006}}{\pi} \qquad \mathrm{(B)}\ \frac {1003}4 \qquad \mathrm{(C)}\ 8\sqrt {1003} \qquad \mathrm{(D)}\ 6\sqrt {2006} \qquad \mathrm{(E)}\ \frac {32\sqrt {1003}}\pi$

 \textbf{Answer: }C



\textbf{Solution}


                    \begin{figure}[H]
                        \centering
                        \includegraphics[width=0.5\textwidth]{Solution/AMC_12_12B_2006_21_Solution_Figure_1.png}
                    \end{figure}
                    

 Let the rectangle have side lengths $l$ and $w$.  Let the axis of the ellipse on which the foci lie have length $2a$, and let the other axis have length $2b$.  We have
 \[lw=ab=2006\]From the definition of an ellipse, $l+w=2a\Longrightarrow \frac{l+w}{2}=a$.  Also, the diagonal of the rectangle has length $\sqrt{l^2+w^2}$.  Comparing the lengths of the axes and the distance from the foci to the center, we have
 \[a^2=\frac{l^2+w^2}{4}+b^2\Longrightarrow \frac{l^2+2lw+w^2}{4}=\frac{l^2+w^2}{4}+b^2\Longrightarrow \frac{lw}{2}=b^2\Longrightarrow b=\sqrt{1003}\]Since $ab=2006$, we now know $a\sqrt{1003}=2006\Longrightarrow a=2\sqrt{1003}$ and because $a=\frac{l+w}{2}$, or one-fourth of the rectangle's perimeter, we multiply by four to get an answer of $\boxed{8\sqrt{1003}}$.

 


	\item (\textit{AMC 12 12B 2006}) Suppose $a$, $b$ and $c$ are positive integers with $a+b+c=2006$, and $a!b!c!=m\cdot 10^n$, where $m$ and $n$ are integers and $m$ is not divisible by $10$. What is the smallest possible value of $n$?


 $\mathrm{(A)}\ 489 \qquad \mathrm{(B)}\ 492  \qquad \mathrm{(C)}\ 495 \qquad \mathrm{(D)}\ 498 \qquad \mathrm{(E)}\ 501$

 \textbf{Answer: }B



\textbf{Solution 1}

The power of $10$ for any factorial is given by the well-known algorithm 
 \[\left\lfloor \frac n{5}\right\rfloor + \left\lfloor \frac n{25}\right\rfloor + \left\lfloor \frac n{125}\right\rfloor + \cdots\]It is rational to guess numbers right before powers of $5$ because we won't have any extra numbers from higher powers of $5$. As we list out the powers of 5, it is clear that $5^{4}=625$ is less than 2006 and  $5^{5}=3125$ is greater. Therefore, set $a$ and $b$ to be 624. Thus, c is $2006-(624\cdot 2)=758$. Applying the algorithm, we see that our answer is $152+152+188= \boxed{492}$.

 



\textbf{Solution 2}

Clearly, the power of $2$ that divides $n!$ is larger or equal than the power of $5$ which dividesit. Hence we are trying to minimize the power of $5$ that will divide 

 


	\item (\textit{AMC 12 12B 2006}) Isosceles $\triangle ABC$ has a right angle at $C$.  Point $P$ is inside $\triangle ABC$, such that $PA=11$, $PB=7$, and $PC=6$. Legs $\overline{AC}$ and $\overline{BC}$ have length $s=\sqrt{a+b\sqrt{2}}$, where $a$ and $b$ are positive integers.  What is $a+b$?


 
                    \begin{figure}[H]
                        \centering
                        \includegraphics[width=0.5\textwidth]{Problem/AMC_12_12B_2006_23_Problem_Figure_1.png}
                    \end{figure}
                    
$\mathrm{(A)}\ 85 \qquad \mathrm{(B)}\ 91 \qquad \mathrm{(C)}\ 108 \qquad \mathrm{(D)}\ 121 \qquad \mathrm{(E)}\ 127$

 \textbf{Answer: }E



\textbf{Solution}


                    \begin{figure}[H]
                        \centering
                        \includegraphics[width=0.5\textwidth]{Solution/AMC_12_12B_2006_23_Solution_Figure_1.png}
                    \end{figure}
                    Using the Law of Cosines on $\triangle PBC$, we have:

 
 \begin{align*} PB^2&=BC^2+PC^2-2\cdot BC\cdot PC\cdot \cos(\alpha) \Rightarrow 49 = 36 + s^2 - 12s\cos(\alpha) \Rightarrow \cos(\alpha) = \dfrac{s^2-13}{12s}. \end{align*}


 Using the Law of Cosines on $\triangle PAC$, we have:
 \begin{align*} PA^2&=AC^2+PC^2-2\cdot AC\cdot PC\cdot \cos(90^\circ-\alpha) \Rightarrow 121 = 36 + s^2 - 12s\sin(\alpha) \Rightarrow \sin(\alpha) = \dfrac{s^2-85}{12s}. \end{align*}


 Now we use $\sin^2(\alpha) + \cos^2(\alpha) = 1$.
 \begin{align*} \sin^2(\alpha)+\cos^2(\alpha) = 1 &\Rightarrow \frac{s^4-26s^2+169}{144s^2} + \frac{s^4-170s^2+7225}{144s^2} = 1 \\ &\Rightarrow s^4-170s^2+3697 = 0 \\ &\Rightarrow s^2 = \dfrac{170 \pm 84\sqrt{2}}{2} = 85 \pm 42\sqrt2  \end{align*}


 Note that we know that we want the solution with $s^2 > 85$ since we know that $\sin(\alpha) > 0$.  Thus, $a+b=85+42=\boxed{127}$.

 



\textbf{Solution 2}

Rotate triangle $PAC$ 90 degrees counterclockwise about $C$ so that the image of $A$ rests on $B$. Now let the image of $P$ be $\widetilde{P}$. 
                    \begin{figure}[H]
                        \centering
                        \includegraphics[width=0.5\textwidth]{Solution/AMC_12_12B_2006_23_Solution_Figure_2.png}
                    \end{figure}
                    Note that $\widetilde{P}C=6$, meaning triangle $PC\widetilde{P}$ is right isosceles, and $\angle P\widetilde{P}C=45^\circ$. Then $P\widetilde{P}=6\sqrt{2}$. Now because $PB=7$ and $\widetilde{P}B=11$, we observe that $\angle \widetilde{P}PB=90^\circ$, by the Pythagorean Theorem on $\widetilde{P}PB$. Now we have that $\angle APC=\angle B\widetilde{P}C=\angle B\widetilde{P}P + \angle P\widetilde{P}C$. So we take the cosine of the second equality, using that fact that $\angle P\widetilde{P}C=45^\circ$, to get $\cos(B\widetilde{P}C)=\frac{6\sqrt{2}-7}{11\sqrt{2}}$. Finally, we use the fact that $\cos(B\widetilde{P}C)=\cos(CPA)$ and use the Law of Cosines on triangle $CPA$ to arrive at the value of $s^2$.

 Or notice that since $\angle \widetilde{P}PB=90^\circ$ and $\angle P\widetilde{P}C=45^\circ$, we have $\angle BPC=135^\circ$, and Law of Cosines on triangle $BPC$ gives the value of $s^2$.

 



\textbf{Solution 3 (coordinate bash)}

Let point $P$ have coordinates $(x,y)$ and $C$ have coordinates $(0,0).$ Then, $A$ has $(s,0)$ and $B$ has $(0,s)$. 

 By distance formula, we have 
 \[x^2+y^2=36 \tag{1}.\]
 \[x^2+(y-s)^2=49 \tag{2}.\]
 \[(x-s)^2+y^2=121 \tag{3}\] 

 Expanding $(2)$ and $(3)$ gives 
 \[x^2+y^2-2ys+s^2=49,\] and 
 \[x^2+y^2-2sx+s^2=121,\] respectively. Then, using equation $(1)$, we have that 
 \[-2ys+s^2=49-36=13,\] and 
 \[-2sx+s^2=121-36=85.\] 

 Then, solving for $x$ and $y$ gives
 \[x=\frac{85-s^2}{-2s}=\frac{s^2-85}{2s},\] and
 \[y=\frac{13-s^2}{-2s}=\frac{s^2-13}{2s}.\] Plugging these values of $x$ and $y$ into equation $(1)$ yields 
 \[\left(\frac{s^2-85}{2s}\right)^2+\left(\frac{s^2-13}{2s}\right)^2=36.\]

 We multiply both sides by $4s^2$ and expand, yielding the equation 
 \[s^4-170s^2+7225+s^4-26s^2+169=144s^2.\] Simplifying gives the equation 
 \[2s^4-340s^2+7394=0.\]  Solving this quadratic gives $s^2=85\pm 42\sqrt{2}.$ Now, if this were the actual test, we stop here, noting that the question tells us $a$ and $b$ are positive, so $s^2$ must be $85+42\sqrt{2}$, and our solution is $127$.

 However, here is why $s^2$ cannot be $85-42\sqrt{2}$:

 If $s^2=85-42\sqrt{2}$, using $1.4$ as an approximation for $\sqrt{2}$, $s^2\approx 85-42\cdot 1.4=85-58.8=26.2$, and $s$ is slightly greater than $5$. Also note that this implies that $AB\approx 7$. 

 Note that at least one of $\angle BPA$, $\angle APC$ and $\angle BPC$ must be obtuse, since they sum to $360^{\circ}$. Then, note the well known fact that if $\angle A$ is the largest angle in $\Delta ABC$, $BC$ must be the largest side. However, combined with the first fact, implies that either $AB$ is the largest side of $\Delta ABP$, $BC$ is the largest side of $\Delta BPC$, or $AC$ is the largest side of $\Delta APC$. By our approximations, this cannot possibly be true, even if we are generous with our margin of error, so $s^2$ cannot equal $85-42\sqrt{2}$, and $s^2=85+42\sqrt{2}$. 

 The answer is $85+42=\boxed{127}$



	\item (\textit{AMC 12 12B 2006}) Let $S$ be the set of all points $(x,y)$ in the coordinate plane such that $0\leq x\leq \frac\pi 2$ and $0\leq y\leq \frac\pi 2$. What is the area of the subset of $S$ for which $\sin^2 x - \sin x\sin y + \sin^2 y\le \frac 34$?


 $\mathrm{(A)}\ \frac {\pi^2}9 \qquad \mathrm{(B)}\ \frac {\pi^2}8 \qquad \mathrm{(C)}\ \frac {\pi^2}6 \qquad \mathrm{(D)}\ \frac {3\pi^2}{16}  \qquad \mathrm{(E)}\ \frac {2\pi^2}9$

 \textbf{Answer: }C



\textbf{Solution}

We start out by solving the equality first.
 \begin{align*} \sin^2x - \sin x \sin y + \sin^2y &= \frac34 \\ \sin x &= \frac{\sin y \pm \sqrt{\sin^2 y - 4 ( \sin^2y - \frac34 ) }}{2} \\  &= \frac{\sin y \pm \sqrt{3 - 3 \sin^2 y }}{2} \\  &= \frac{\sin y \pm \sqrt{3 \cos^2 y }}{2} \\  &= \frac12 \sin y \pm \frac{\sqrt3}{2} \cos y \\ \sin x &= \sin (y \pm \frac{\pi}{3}) \end{align*}
We end up with three lines that matter: $x = y + \frac\pi3$, $x = y - \frac\pi3$, and $x = \pi - (y + \frac\pi3) = \frac{2\pi}{3} - y$.  We plot these lines below.
                    \begin{figure}[H]
                        \centering
                        \includegraphics[width=0.5\textwidth]{Solution/AMC_12_12B_2006_24_Solution_Figure_1.png}
                    \end{figure}
                    Note that by testing the point $(\pi/6,\pi/6)$, we can see that we want the area of the \href{https://artofproblemsolving.com/wiki/index.php?title=Pentagon}{pentagon}.  We can calculate that by calculating the area of the square and then subtracting the area of the 3 triangles. (Note we could also do this by adding the areas of the isosceles triangle in the bottom left corner and the rectangle with the previous triangle's hypotenuse as the longer side.)
 \begin{align*} A &= \left(\frac{\pi}{2}\right)^2 - 2 \cdot \frac12 \cdot \left(\frac{\pi}{6}\right)^2 - \frac12 \cdot \left(\frac{\pi}{3}\right)^2 \\ &= \pi^2 \left ( \frac14 - \frac1{36} - \frac1{18}\right ) \\ &= \pi^2 \left ( \frac{9-1-2}{36} \right ) = \boxed{\text{(C)}\ \frac{\pi^2}{6}} \end{align*}





\textbf{Solution 2}

We can write the given equation as 
 \[\sin^3x + \sin^3y \le \frac{3}{4}(\sin x + \sin y).\] Note that when $x = 0$, we have $\sin y \le \frac{\sqrt{3}}{2}$ which implies $y \le \frac{\pi}{3}$. Similary we have $x \le \frac{\pi}{3}$ when $y = 0$. Then we see what happens at $x = \frac{\pi}{2}$. Clearly at $x = \frac{\pi}{2}$, we see that $y \le \frac{\pi}{6}$. By symmetry we have $x \le \frac{\pi}{6}$ when $y = \frac{\pi}{2}$. So we get a graph like 
                    \begin{figure}[H]
                        \centering
                        \includegraphics[width=0.5\textwidth]{Solution/AMC_12_12B_2006_24_Solution_Figure_1.png}
                    \end{figure}
                    Thus the area of what we are interested in is 
 \[\frac{\pi^2}{4} - (\frac{\pi^2}{18} + \frac{\pi^2}{36}) = \frac{\pi^2}{6}.\]~coolmath\_2018

 


	\item (\textit{AMC 12 12B 2006}) A sequence $a_1,a_2,\dots$ of non-negative integers is defined by the rule $a_{n+2}=|a_{n+1}-a_n|$ for $n\geq 1$. If $a_1=999$, $a_2<999$ and $a_{2006}=1$, how many different values of $a_2$ are possible?


 $\mathrm{(A)}\ 165 \qquad \mathrm{(B)}\ 324 \qquad \mathrm{(C)}\ 495 \qquad \mathrm{(D)}\ 499 \qquad \mathrm{(E)}\ 660$

 \textbf{Answer: }B



\textbf{Solution 1}

We say the sequence $(a_n)$ completes at $i$ if $i$ is the minimal positive integer such that $a_i = a_{i + 1} = 1$. Otherwise, we say $(a_n)$ does not complete.

 
Note that if $d = \gcd(999, a_2) \neq 1$, then $d|a_n$ for all $n \geq 1$, and $d$ does not divide $1$, so if $\gcd(999, a_2) \neq 1$, then $(a_n)$ does not complete. (Also, $a_{2006}$ cannot be 1 in this case since $d$ does not divide $1$, so we do not care about these $a_2$ at all.)

 From now on, suppose $\gcd(999, a_2) = 1$.

 
We will now show that $(a_n)$ completes at $i$ for some $i \leq 2006$. We will do this with 3 lemmas.

 
\textbf{Lemma:} If $a_j \neq a_{j + 1}$, and neither value is $0$, then $\max(a_j, a_{j + 1}) > \max(a_{j + 2}, a_{j + 3})$.

 \textbf{Proof:} There are 2 cases to consider.

 If $a_j > a_{j + 1}$, then $a_{j + 2} = a_j - a_{j + 1}$, and $a_{j + 3} = |a_j - 2a_{j + 1}|$. So $a_j > a_{j + 2}$ and $a_j > a_{j + 3}$.

 If $a_j < a_{j + 1}$, then $a_{j + 2} = a_{j + 1} - a_j$, and $a_{j + 3} = a_j$. So $a_{j + 1} > a_{j + 2}$ and $a_{j + 1} > a_{j + 3}$.

 In both cases, $\max(a_j, a_{j + 1}) > \max(a_{j + 2}, a_{j + 3})$, as desired.

 \textbf{Lemma:} If $a_i = a_{i + 1}$, then $a_i = 1$. Moreover, if instead we have $a_i = 0$ for some $i > 2$, then $a_{i - 1} = a_{i - 2} = 1$.

 \textbf{Proof:} By the way $(a_n)$ is constructed in the problem statement, having two equal consecutive terms $a_i = a_{i + 1}$ implies that $a_i$ divides every term in the sequence. So $a_i | 999$ and $a_i | a_2$, so $a_i | \gcd(999, a_2) = 1$, so $a_i = 1$. For the proof of the second result, note that if $a_i = 0$, then $a_{i - 1} = a_{i - 2}$, so by the first result we just proved, $a_{i - 2} = a_{i - 1} = 1$.

 \textbf{Lemma:} $(a_n)$ completes at $i$ for some $i \leq 2000$.

 \textbf{Proof:} Suppose $(a_n)$ completed at some $i > 2000$ or not at all. Then by the second lemma and the fact that neither $999$ nor $a_2$ are $0$, none of the pairs $(a_1, a_2), ..., (a_{1999}, a_{2000})$ can have a $0$ or be equal to $(1, 1)$. So the first lemma implies
 \[\max(a_1, a_2) > \max(a_3, a_4) > \cdots > \max(a_{1999}, a_{2000})  > 0,\]so $999 = \max(a_1, a_2) \geq 1000$, a contradiction. Hence $(a_n)$ completes at $i$ for some $i \leq 2000$.

 
Now we're ready to find exactly which values of $a_2$ we want to count.

 Let's keep in mind that $2006 \equiv 2 \pmod 3$ and that $a_1 = 999$ is odd. We have two cases to consider.

 
\textbf{Case 1:} If $a_2$ is odd, then $a_3$ is even, so $a_4$ is odd, so $a_5$ is odd, so $a_6$ is even, and this pattern must repeat every three terms because of the recursive definition of $(a_n)$, so the terms of $(a_n)$ reduced modulo 2 are
 \[1, 1, 0, 1, 1, 0, ...,\] so $a_{2006}$ is odd and hence $1$ (since if $(a_n)$ completes at $i$, then $a_k$ must be $0$ or $1$ for all $k \geq i$).

 \textbf{Case 2:} If $a_2$ is even, then $a_3$ is odd, so $a_4$ is odd, so $a_5$ is even, so $a_6$ is odd, and this pattern must repeat every three terms, so the terms of $(a_n)$ reduced modulo 2 are 
 \[1, 0, 1, 1, 0, 1, ...,\]so $a_{2006}$ is even, and hence $0$.

 
We have found that $a_{2006} = 1$ is true precisely when $\gcd(999, a_2) = 1$ and $a_2$ is odd. This tells us what we need to count.

 
There are $\phi(999) = 648$ numbers less than $999$ and relatively prime to it ($\phi$ is the Euler totient function). We want to count how many of these are odd. Note that
 \[t \mapsto 999 - t\]is a 1-1 correspondence between the odd and even numbers less than and relatively prime to $999$. So our final answer is $648/2 = 324$, or $\boxed{\text{B}}$.

 



\textbf{Solution 2 (estimation)}

As in Solution 1, we show that $a_2$ must be odd. We notice that, after experimentation, all numbers eventually end up in a sequence $1,1,0,1,1,0,\ldots$, which points to the assumption that every $2$ out of $3$ values for $a_2$ work due to residue classes. Combining this with an estimated 50-50 chance that these values are either odd or even, since $999$ values are nonnegative integers less than $999$, we can estimate that there are about $999\cdot\frac{2}{3}\cdot\frac{1}{2}=333$ different values of $a_2$ that work. The closest answer choice by far is $\boxed{\textbf{(B) }324}$.

 



\textbf{Solution 3 (PIE)}

Notice that if $a_2$ is not relatively prime to $999 = 3^3\cdot 37$, then all elements in the sequence will be a multiple of 3 or 37, i.e. it cannot be 1. In addition, if $a_2$ is even, then $a_{2006}$ must also be even, so $a_{2006}\neq 1$. For all other possible $a_2$ (those relatively prime to 2, 3, and 37), the sequence will become $...,1,1,0,1,1,0,...$ where $a_{2006}=1$. From PIE, the complement is
 \[\left(\left\lfloor\frac{999}{2}\right\rfloor+1\right)+\left(\left\lfloor\frac{999}{3}\right\rfloor+1\right)+\left(\left\lfloor\frac{999}{37}\right\rfloor+1\right)-\left(\left\lfloor\frac{999}{2\cdot3}\right\rfloor+1\right)-\left(\left\lfloor\frac{999}{2\cdot 37}\right\rfloor+1\right)-\left(\left\lfloor\frac{999}{3\cdot37}\right\rfloor+1\right)+\left(\left\lfloor\frac{999}{2\cdot3\cdot37}\right\rfloor+1\right)=676\]Note the +1 in each term because we have to include 0.So the answer is $1000-676=\boxed{\textbf{(B) }324}$.

 


\end{enumerate}